\documentclass[11pt]{article}
\usepackage[T1]{fontenc}
\usepackage[margin=1in]{geometry}
\usepackage{amsmath,amssymb,amsthm,booktabs,array}
\usepackage[hidelinks]{hyperref}
\newtheorem{theorem}{Theorem}
\newtheorem{lemma}[theorem]{Lemma}
\newtheorem{corollary}[theorem]{Corollary}
\newcommand{\psd}{\succeq0}
\newcommand{\HH}{\mathcal H}
\newcommand{\Mult}{\operatorname{Mult}}
\newcommand{\diag}{\operatorname{diag}}
\title{A three-point negative answer to\newline
Luo--McCullough--Tsikalas's\newline pseudo-Pick quotient question}
\author{Haobo Ma\thanks{ChronoAI Pte Ltd, Singapore.}\and
Wenlin Zhang\thanks{National University of Singapore; ChronoAI Pte Ltd, Singapore.}}
\date{18 November 2026\\[3pt]\small CC BY 4.0}

\begin{document}
\maketitle
\begin{abstract}
We answer Question 9.3 of Luo, McCullough and Tsikalas,
\emph{Beurling Criteria for Reproducing Kernels}, arXiv:2607.25088v2
(2026), page 45, negatively on a three-point set: a rank-two source kernel
forces constant full interpolation data, while positive retraction quotients
extend every proper datum at every matrix size.
The source kernel is non-vanishing, the target kernel is strictly positive
definite, and their entrywise quotient has quadratic form $-10/33$ at
$(-2,1,1)$. Three points and source rank two are minimal. A diagonal
perturbation gives a one-parameter family with the same negative answer.
The family is neither complete Pick, Beurling--Lax--Halmos, complete
pseudo-Leech, nor weakly Agler-decomposable.
\end{abstract}

\section{The question and the construction}
A complete pseudo-Pick pair need not have a positive quotient, even on
three points. The property in \cite[Definition 9.2]{LMT} concerns interpolation
on the target space itself. This differs from the complete Pick property
for a pair, which interpolates multipliers from the source space to the
target space \cite[Definition 3.25]{LMT}. The familiar one-kernel matrix
interpolation property is described in \cite[Section 0]{AM}; keeping the
two multiplier spaces distinct is essential here. For the mixed-space
property, McCullough and Tsikalas characterize diagonal holomorphic
pairs by Shimorin factorizations \cite[Theorem 1.4]{MT}.
The present construction concerns the same-target-space property of
\cite[Definition 9.2]{LMT} on a finite set, with a singular source kernel.
The separation from mixed-space complete Pick alone has an earlier
antecedent: \cite[Proposition 7.4]{MT}, already in its 2024 preprint,
has a complete Pick target and positive quotient $1+u+4u^2$, where
$u=z\overline w$. Schur multiplication transfers admissible source
data to target Pick data, so that example is complete pseudo-Pick.
Our distinction is the negative quotient, with minimal cardinality
and source rank, not the existence of a separation by itself.

A scalar kernel $h$ on a set $X$ is positive if every finite matrix
$[h(z_i,z_j)]$ is positive semidefinite; repetitions among the $z_i$ are
allowed. Its reproducing kernel Hilbert space is denoted by $\HH_h$.
For a positive integer $m$, let $M_m(\mathbb C)$ be the $m\times m$
complex matrices, $I_m$ their identity, and $W^*$ the conjugate transpose
of $W$. A function $\Phi:X\to M_m(\mathbb C)$ is a contractive
multiplier on $\HH_h\otimes\mathbb C^m$ precisely when
\begin{equation}\label{eq:mult}
  [\,h(i,j)(I_m-\Phi(i)\Phi(j)^*)\,]_{i,j\in X}\psd.
\end{equation}
For finite $X$, this is the ordinary block-matrix criterion; in general
it is a positive-kernel criterion \cite[Section 3, p.~15]{LMT}.
Indeed, on the span of kernel vectors, the adjoint of multiplication must
send $h_j\otimes u$ to $h_j\otimes\Phi(j)^*u$.
The norm inequality for this assignment is exactly \eqref{eq:mult}; it
also makes the assignment well defined, and its bounded extension gives
the multiplier.

A pair $(k,\ell)$ is \emph{complete pseudo-Pick} if for all positive
integers $n,m$, all lists $z_1,\ldots,z_n\in X$ and all
$W_1,\ldots,W_n\in M_m(\mathbb C)$, positivity of
\begin{equation}\label{eq:pick}
  P=[\,k(z_i,z_j)(I_m-W_iW_j^*)\,]_{i,j=1}^n
\end{equation}
implies that a contractive multiplier on $\HH_\ell\otimes\mathbb C^m$
has values $\Phi(z_i)=W_i$. Question 9.3 asks whether this property and
the non-vanishing of $k$ imply positivity of the entrywise quotient
$\ell/k$. In particular, the question imposes no strict positive
definiteness hypothesis on $k$.

Take $X=\{0,1,2\}$, put $x=(1/10,1,2)$, and define
\begin{equation}\label{eq:kernels}
 k(i,j)=1+x_i x_j,\qquad
 \ell(i,j)=1+\delta_{ij}x_i^2,
\end{equation}
where $\delta_{ij}=1$ for $i=j$ and $0$ otherwise. Write
$\mathbf1=(1,1,1)^T$. Their matrices are
\[
 K=\mathbf1\mathbf1^*+xx^T
 =\begin{pmatrix}101/100&11/10&6/5\\11/10&2&3\\6/5&3&5\end{pmatrix},
 \qquad
 E=\mathbf1\mathbf1^*+\diag(1/100,1,4).
\]
Thus $K\psd$ has rank two, all its entries are positive, and $E$ is
positive definite.

\begin{theorem}\label{thm:main}
The kernels in \eqref{eq:kernels} form a complete pseudo-Pick pair.
Their quotient is not positive semidefinite. No pair on a set with at
most two points can have these two properties when its source kernel
is non-vanishing. Source rank two is minimal among finite-set examples.
\end{theorem}

The quotient and a negative vector are explicit:
\begin{equation}\label{eq:negative}
 Q=\left[\frac{\ell(i,j)}{k(i,j)}\right]
 =\begin{pmatrix}1&10/11&5/6\\10/11&1&1/3\\5/6&1/3&1\end{pmatrix},
 \qquad
 (-2,1,1)Q(-2,1,1)^T=-\frac{10}{33}.
\end{equation}
It remains to prove interpolation at every matrix size and the
minimality assertion.

\section{Interpolation for the example}
\begin{lemma}\label{lem:relations}
Suppose the Pick matrix \eqref{eq:pick} is positive for the source kernel
in \eqref{eq:kernels}. If real coefficients $a_1,\ldots,a_n$ satisfy
$\sum_i a_i=\sum_i a_i x_{z_i}=0$, then
\[
 \sum_i a_iW_i=0,\qquad \sum_i a_i x_{z_i}W_i=0.
\]
Consequently, repeated nodes have equal matrix values, and a datum
containing all three distinct nodes has a single constant matrix value.
\end{lemma}
\begin{proof}
Compress $P$ by the block row $[a_1I_m\ \cdots\ a_nI_m]$.
The two scalar relations cancel its positive terms and give
\[
  -AA^*-BB^*\psd,\qquad
  A=\sum_i a_iW_i,\quad B=\sum_i a_ix_{z_i}W_i.
\]
For every vector $u$, this says
$-\|A^*u\|^2-\|B^*u\|^2\geq0$, so $A=B=0$.
For equal nodes, coefficients $1,-1$ at the two occurrences and zero
elsewhere give equality of the matrix values.

If all three nodes occur, select one occurrence of each. Positivity
passes to this principal block matrix, and the coefficients
$(-10,19,-9)$ have zero sum and zero coordinate-weighted sum.
The two resulting relations are
\[
 -10W_0+19W_1-9W_2=0,\qquad -W_0+19W_1-18W_2=0.
\]
Subtracting gives $W_0=W_2$, and substitution gives $W_1=W_0$.
The repeated-node conclusion covers the rest of the list.
\end{proof}

For $a\in X$, define a retraction $r_a:X\to X\setminus\{a\}$ that fixes
the other two points by the tuples
\[
\begin{aligned}
(r_0(0),r_0(1),r_0(2))&=(1,1,2),\\
(r_1(0),r_1(1),r_1(2))&=(0,0,2),\\
(r_2(0),r_2(1),r_2(2))&=(0,1,0).
\end{aligned}
\]
Set $R_a(i,j)=\ell(i,j)/k(r_a(i),r_a(j))$.

\begin{lemma}\label{lem:retract}
Every $R_a$ is positive definite.
\end{lemma}
\begin{proof}
The factorizations $R_a=L_aD_aL_a^T$ are given in the following table.
Here $L_a$ is unit lower triangular, with its three entries below the
diagonal listed in the order $(L_{10},L_{20},L_{21})$, and the final
column gives the diagonal entries of $D_a$.
\begin{center}
\renewcommand{\arraystretch}{1.65}
\begin{tabular}{c c c}
\toprule
$a$ & Entries of $L_a$ & Entries of $D_a$\\
\midrule
$0$ & $\left(\frac{100}{101},\frac{200}{303},\frac{1}{153}\right)$
    & $\left(\frac{101}{200},\frac{51}{101},\frac{358}{459}\right)$\\
$1$ & $\left(\frac{100}{101},\frac56,\frac{101}{12240}\right)$
    & $\left(1,\frac{10200}{10201},\frac{4487}{14688}\right)$\\
$2$ & $\left(\frac{10}{11},\frac{100}{101},\frac{110}{2121}\right)$
    & $\left(1,\frac{21}{121},\frac{850400}{214221}\right)$\\
\bottomrule
\end{tabular}
\end{center}
Expanding the matrix products verifies each identity. All entries of $D_a$
are positive and each $L_a$ is invertible.
\end{proof}

We also use the block Schur product fact: if $[A_{ij}]\psd$ is a block
matrix and $[s_{ij}]\psd$ is a scalar matrix, then
$[s_{ij}A_{ij}]\psd$. To see this, factor
$s_{ij}=\sum_\nu c_{i\nu}\overline{c_{j\nu}}$ and write the result as
$\sum_\nu C_\nu[A_{ij}]C_\nu^*$, where
$C_\nu=\diag(c_{i\nu}I_m)$.

\begin{proof}[Proof of interpolation in Theorem~\ref{thm:main}]
Fix any positive $n,m$ and any positive Pick datum. A diagonal block of
$P$ shows $I_m-W_iW_i^*\psd$, since $k(z_i,z_i)>0$.
By Lemma~\ref{lem:relations}, the values at repeated nodes agree.
If one distinct node occurs, extend its value $U$ constantly to $X$.
If all three nodes occur, Lemma~\ref{lem:relations} again gives one
value $U$, and use the same extension. In both cases,
$[\ell(i,j)(I_m-UU^*)]\psd$ because it is a tensor product of two
positive matrices.

If exactly two distinct nodes occur, let $a$ be the missing node and
denote their values by $V_b$ for $b\ne a$. Define
$\Phi(i)=V_{r_a(i)}$. This agrees with every prescribed value. The
matrix
\[
 A_{ij}=k(r_a(i),r_a(j))
       (I_m-V_{r_a(i)}V_{r_a(j)}^*)
\]
is positive: it is obtained by repeating rows and columns of the positive
two-node Pick matrix, equivalently by a matrix congruence. Moreover,
\[
 \ell(i,j)(I_m-\Phi(i)\Phi(j)^*)=R_a(i,j)A_{ij}.
\]
Lemma~\ref{lem:retract} and the block Schur product fact prove positivity
of the target matrix. Criterion \eqref{eq:mult} gives the required
contractive multiplier. The cases exhaust all finite lists, including
lists with repetitions, for every $m$.
\end{proof}

\section{Minimality and related properties}
\begin{lemma}\label{lem:two}
Let $(k,\ell)$ be complete pseudo-Pick on a set with at most two points,
with $k$ non-vanishing. Then $\ell/k$ is positive.
\end{lemma}
The two-point lower bound is also a consequence of the classical
multiplier-distance inequality \cite[Proposition 4]{ARSW}; the following
argument gives the quotient determinant directly, including degenerate
target kernels.
\begin{proof}
The empty and one-point cases follow from nonnegative diagonal entries.
For two points, write $k_{ij}=k(i,j)$ and $\ell_{ij}=\ell(i,j)$.
Non-vanishing and positivity of $k$ give $k_{00},k_{11}>0$ and
\[
  \rho=\frac{|k_{01}|}{\sqrt{k_{00}k_{11}}}\in(0,1].
\]
Put $t=\sqrt{(1-\rho)/(1+\rho)}$. The scalar datum $W_0=t$, $W_1=-t$
has source Pick matrix
\[
 \begin{pmatrix}(1-t^2)k_{00}&(1+t^2)k_{01}\\
 (1+t^2)k_{10}&(1-t^2)k_{11}\end{pmatrix}\psd:
\]
its diagonal is nonnegative and its determinant is zero.
Pseudo-Pick interpolation and \eqref{eq:mult} give
\[
 (1-t^2)^2\ell_{00}\ell_{11}
 \geq (1+t^2)^2|\ell_{01}|^2.
\]
Since $(1-t^2)/(1+t^2)=\rho$, this is
$|\ell_{01}|^2\leq\rho^2\ell_{00}\ell_{11}$, which makes the
determinant of $[\ell_{ij}/k_{ij}]$ nonnegative. Its diagonal is
nonnegative as well. This also covers a zero diagonal entry of $\ell$.
\end{proof}
Together with \eqref{eq:negative}, this proves the point-minimality
assertion of Theorem~\ref{thm:main}.

For its rank assertion, a non-vanishing rank-one positive kernel on a
finite set has the form $k(i,j)=u_i\overline{u_j}$ with every $u_i\ne0$.
For any positive target matrix $E$ the quotient is
$D^{-1}E(D^{-1})^*$, where $D=\diag(u_i)$, and is therefore positive.
Rank zero is excluded by non-vanishing.

A Beurling--Lax--Halmos (BLH) pair requires every
$\Mult(\HH_k)$-invariant subspace of $\HH_\ell\otimes\mathcal E$
to be the range of a partially isometric multiplier from some
$\HH_k\otimes\mathcal F$ \cite[Definition 1.2]{LMT}.

\begin{corollary}\label{cor:family}
For $0\leq t<5/99$, define
\[
 \ell_t(i,j)=\ell(i,j)+t\delta_{ij}k(i,i).
\]
Then $(k,\ell_t)$ is complete pseudo-Pick, its target is positive
definite, and its quotient is not positive. Each pair is not BLH,
not complete pseudo-Leech, and admits neither a weak nor a strong
Agler decomposition.
\end{corollary}
\begin{proof}
The target matrix is $E+t\diag(k(0,0),k(1,1),k(2,2))$, hence is
positive definite. The retraction quotients become
\[
 R_{a,t}=R_a+t\diag\left(
 \frac{k(i,i)}{k(r_a(i),r_a(i))}\right)_{i\in X},
\]
which remain positive definite for every $t\geq0$. The proof in
Section~2, including repeated nodes and all matrix sizes, is unchanged.
The entrywise quotient is $Q+tI_3$, so
\[
 (-2,1,1)(Q+tI_3)(-2,1,1)^T=-\frac{10}{33}+6t<0
 \qquad (0\leq t<5/99).
\]
This interval is sufficient; no optimal endpoint is asserted.

Every $\HH_{\ell_t}$ consists of all functions on $X$ and so is
$\Mult(\HH_k)$-invariant. If the pair were BLH, applying
\cite[Proposition 3.1]{LMT} to this entire target space would give a
positive kernel $g$ with $\ell_t=kg$. Non-vanishing of $k$ would force
the indefinite quotient $\ell_t/k$ to be positive, a contradiction.

For the pseudo-Leech assertion, recall that complete pseudo-Leech means
that every positive kernel of the form
$[k(i,j)(A_iA_j^*-B_iB_j^*)]$ factors as
$B_i=A_i\Phi(i)$ through a contractive target-space multiplier, for
arbitrary finite rectangular matrix data
\cite[Definition 4.1, pp.~23--24]{LMT}.
The argument of \cite[Proposition 4.3]{LMT}, applied to the entire target
space, shows directly that this property requires a positive quotient.
Indeed, on any finite set, the Hermitian matrix $[1/k(i,j)]$ has a
decomposition $AA^*-BB^*$, by its positive and negative spectral parts;
pad both factors with zero columns to the same width $d$.
Its rows $A_i,B_i\in\mathbb C^{1\times d}$ give the positive input
$[k(i,j)(A_iA_j^*-B_iB_j^*)]=\mathbf1\mathbf1^*$.
A pseudo-Leech interpolant would imply positivity of
$[\ell_t(i,j)(I_d-\Phi(i)\Phi(j)^*)]$.
Compressing this block matrix by the block-diagonal row map with rows
$A_i$ gives precisely $[\ell_t(i,j)/k(i,j)]$, a contradiction.
This deduction uses no invariant-subspace comparison hypothesis.

Finally, a weak Agler decomposition is a pointwise limit
\[
 Q_n(i)(I-\Phi_n(i)\Phi_n(j)^*)Q_n(j)^*
       \longrightarrow 1/k(i,j),
\]
where $\Phi_n$ are contractive target-space multipliers and $Q_n(i)$
are finite row vectors \cite[Definition 3.13]{LMT}.
Multiplication by $\ell_t(i,j)$ makes every approximant positive,
by the same block compression. Its limit would be $\ell_t/k$,
which is not positive. This is the necessary condition already noted
in \cite[Remark 3.14]{LMT}. A strong Agler decomposition would give
the same positivity directly \cite[Definition 3.17]{LMT}.
\end{proof}

\begin{corollary}\label{cor:separation}
Every pair $(k,\ell_t)$ in Corollary~\ref{cor:family} is neither a
complete Pick pair nor a Beurling--Lax--Halmos pair.
\end{corollary}
\begin{proof}
Shimorin's necessary factorization theorem for mixed complete Pick pairs
\cite[Theorem 1.4]{Shimorin}, also restated as
\cite[Theorem 1.9]{MT}, already implies that every pair in the family
fails complete Pick, because such a pair would have a positive quotient.
The following scalar datum gives an explicit obstruction as well.
For the complete Pick property for a pair, the input matrix is
$[\ell(z_i,z_j)I_m-k(z_i,z_j)W_iW_j^*]$, and the interpolant is a
contractive multiplier from $\HH_k\otimes\mathbb C^m$ to
$\HH_\ell\otimes\mathbb C^m$ \cite[Definition 3.25]{LMT}.
The single scalar datum $\Phi(0)=1$ is admissible because
$\ell(0,0)=k(0,0)$. If such a mixed-space multiplier existed, its
positive kernel
$B_{ij}=\ell(i,j)-k(i,j)\Phi(i)\overline{\Phi(j)}$ would satisfy
$B_{00}=0$. Every positive matrix with a zero diagonal entry has zero
in the corresponding row: apply its $2\times2$ principal minors.
It follows that
$\Phi=(1,10/11,5/6)$. But then
\[
 [B_{ij}]_{i,j=1,2}=
 \begin{pmatrix}42/121&-14/11\\-14/11&55/36\end{pmatrix},
 \qquad \det[B_{ij}]_{i,j=1,2}=-\frac{791}{726}<0,
\]
a contradiction for $t=0$.

For general $t$ in the indicated interval, put $s=1+t$,
$a=10/11$ and $b=5/6$. The admissible single datum
$\Phi(0)=\sqrt{s}$ forces
$\Phi=(\sqrt{s},a/\sqrt{s},b/\sqrt{s})$ by the same zero-row argument.
The determinant of the remaining $2\times2$ block is $F(s)/s^2$, where
\[
 F(s)=10(s^2-a^2)(s^2-b^2)-(s-3ab)^2.
\]
Set $C=10(a^2+b^2)+1=35303/2178<20$.
For $s\geq1$, $F'(s)=40s^3-2Cs+6ab>0$.
Since $s<104/99$ and
$F(104/99)=-34591631/96059601<0$, the required determinant is negative.

The BLH assertion for every $t$ is part of Corollary~\ref{cor:family}.
\end{proof}

The discussion preceding Problem 9.4 in \cite[p.~45]{LMT} already
separates pseudo-Pick from complete Pick in a diagonal holomorphic
setting. Corollary~\ref{cor:separation} supplies a finite-set
separation. The paragraph preceding Question 9.3 considers a route from
pseudo-Pick to BLH; it does not assert that implication as a theorem.
Corollary~\ref{cor:family} proves its failure on general sets. The
converse of the stated pseudo-Leech implication also fails. These are
obstructions, not refutations of assertions made in the source.

The source multiplier algebra in this example consists of constants.
Indeed, $\HH_k=\operatorname{span}\{\mathbf1,x\}$ as a vector space.
If $\varphi$ is a multiplier, $\varphi=a+bx$ and
$\varphi x=c+dx$ for scalars $a,b,c,d$. Thus
$bx_i^2+(a-d)x_i-c=0$ at three distinct values $x_i$, so the polynomial
vanishes identically and $b=c=0$, $a=d$. This also explains why the
example does not belong to the domain-based regular setting of the
source's Questions 9.1, 9.5 and 9.6. In particular, failure of BLH here
does not answer the BLH-constrained existence question 9.5.
The construction makes no assertion about strictly positive definite
source kernels or holomorphic-domain variants.
For the following comparison, normalization at $o$ means $h(z,o)=1$
for all $z$, for each of the two kernels; irreducibility means
non-vanishing and linear independence of distinct kernel vectors.
Hartz's multiplier-inclusion criterion \cite[Corollary 3.5]{Hartz}
gives an affirmative subclass: if $k$ is an irreducible complete
Pick kernel and $k,\ell$ are normalized at the same point, complete
pseudo-Pick implies $\ell/k\psd$. Indeed, test every finite restriction
of a matrix-valued contractive source multiplier; pseudo-Pick gives
positivity of its target multiplier kernel on each finite set. This
is completely contractive multiplier inclusion, to which Hartz's
criterion applies.

\begin{center}
\small
\begin{tabular}{p{.26\textwidth}p{.65\textwidth}}
\toprule
Source item in \cite{LMT} & Coverage\\
\midrule
Question 9.3, p.~45 & Answered: no. Three points and source rank two
are minimal; Corollary~\ref{cor:family} gives a family.\\
Pseudo-Pick $\Rightarrow$ BLH discussion, p.~45 & Not asserted as a
theorem in the source; fails for \eqref{eq:kernels}.\\
Problem 9.4, p.~45 & Finite pseudo-Pick/non-complete-Pick separation;
earlier separation: McCullough--Tsikalas, Proposition 7.4 (2024);
the posing source also gives a holomorphic separation. Full comparison
not covered.\\
Pseudo-Leech implication, pp.~24, 45 & The forward implication is stated
in the source; its converse is not asserted there and fails
for \eqref{eq:kernels}.\\
Agler decompositions, Definitions 3.13, 3.17 & The family admits neither;
the quotient obstruction is credited to Remark 3.14.\\
Question 9.1, pp.~44--45 & Not covered: domain and dilation hypotheses.\\
Question 9.5, p.~45 & Not covered: a regular BLH pair is required.\\
Question 9.6, p.~45 & Not covered: polynomial row versus matrix norms.\\
\bottomrule
\end{tabular}
\end{center}

\section{Verification scope and supporting material}
The Lean declaration \texttt{result} in the
\href{https://github.com/the-omega-institute/trureturing/blob/8acf7efa41af7fc9ddbbd39571e7ef73462f4509/D5/S3/HilbertOperators/ReproducingKernels/CompletePseudoPickQuotientRefutation.lean}{pinned Lean file}
proves the negation of \texttt{claim}, the universal quotient-positivity
assertion for non-vanishing positive complex kernels with the complete
pseudo-Pick property. Its proof uses the three-point kernels above and
covers arbitrary positive matrix sizes and repeated interpolation
nodes. Lemma~\ref{lem:two} and Corollary~\ref{cor:separation} are written
proofs, as are the rank-one obstruction, the constant-multiplier
description and all assertions of Corollary~\ref{cor:family}.
In particular the pinned declaration concerns $t=0$, not the whole
parameter family. The supporting \texttt{verify.py} and \texttt{checks.json}
give exact rational checks of \eqref{eq:negative}, the three retraction
factorizations, the source relations and the mixed-multiplier obstruction.
They also give the coefficient $6$ and endpoint $5/99$ in the family's
negative quadratic form, a rational interior example, and its positive
retraction factorizations. These are exact identities; the interval and
all-size conclusions have the written proofs above.

\begin{thebibliography}{9}
\bibitem{AM}
J.~Agler and J.~E.~McCarthy,
\emph{Complete Nevanlinna--Pick kernels},
Journal of Functional Analysis \textbf{175} (2000), 111--124.
DOI: \href{https://doi.org/10.1006/jfan.2000.3599}{10.1006/jfan.2000.3599}.
Author manuscript, Section~0:
\url{https://arxiv.org/abs/1610.01997v1}.

\bibitem{LMT}
S.~Luo, S.~McCullough and G.~Tsikalas,
\emph{Beurling Criteria for Reproducing Kernels},
arXiv:2607.25088v2 (14 August 2026).
\url{https://arxiv.org/abs/2607.25088v2}.

\bibitem{MT}
S.~McCullough and G.~Tsikalas,
\emph{The Complete Pick Property for Pairs of Kernels and Shimorin's
Factorization}, Journal of the London Mathematical Society
\textbf{112} (2025), no.~4, e70325.
DOI: \href{https://doi.org/10.1112/jlms.70325}{10.1112/jlms.70325}.
Author manuscript, Theorem~1.4 and Proposition~7.4:
\url{https://arxiv.org/abs/2405.16319v4}.
First preprint: \url{https://arxiv.org/abs/2405.16319v1} (25 May 2024).

\bibitem{ARSW}
N.~Arcozzi, R.~Rochberg, E.~T.~Sawyer and B.~D.~Wick,
\emph{Distance Functions for Reproducing Kernel Hilbert Spaces},
Contemporary Mathematics \textbf{547} (2011), 25--53.
DOI: \href{https://doi.org/10.1090/conm/547/10805}{10.1090/conm/547/10805}.
Author manuscript, Proposition~4, p.~5:
\url{https://arxiv.org/abs/1010.0136v1}.

\bibitem{Shimorin}
S.~Shimorin,
\emph{Commutant lifting and factorization of reproducing kernels},
Journal of Functional Analysis \textbf{224} (2005), 134--159.
DOI: \href{https://doi.org/10.1016/j.jfa.2004.08.011}{10.1016/j.jfa.2004.08.011}.

\bibitem{Hartz}
M.~Hartz,
\emph{On the isomorphism problem for multiplier algebras of
Nevanlinna--Pick spaces}, Canadian Journal of Mathematics
\textbf{69} (2017), 54--106.
DOI: \href{https://doi.org/10.4153/CJM-2015-050-6}{10.4153/CJM-2015-050-6}.
Author manuscript, Corollary~3.5:
\url{https://arxiv.org/abs/1505.05108v1}.
\end{thebibliography}
\end{document}
